\chapter{HYDRODYNAMIC MODELING AND GRAPH CASCADE ENGINE}

\section{Hydrological Formulation and Catchment Inversion}
Urban stormwater accumulation across a municipality is driven by meteorological precipitation falling upon heterogeneous land surfaces. CivicTwin models peak runoff flow rate using the rational method. For each conduit segment $v \in V$, the expected peak surface runoff $Q(v)$ arriving at the conduit during a storm of intensity $i$ (in $\text{mm}/\text{hr}$) is governed by:
\begin{equation}
Q(v) = \frac{C(v) \cdot i \cdot A(v)}{3.6 \times 10^6} \quad [\text{m}^3/\text{s}]
\label{eqn:rational_specific}
\end{equation}
where $C(v) \in [0.1, 0.95]$ represents the localized dimensionless runoff coefficient, $i$ is the uniform rainfall intensity parameter, and $A(v)$ is the tributary catchment surface area (in $\text{m}^2$) contributing drainage directly to conduit $v$. The constant divisor $3.6 \times 10^6$ reconciles unit conversions between millimeters, hours, square meters, and seconds:
\begin{equation}
1~\frac{\text{mm} \cdot \text{m}^2}{\text{hr}} = \frac{10^{-3}~\text{m} \cdot \text{m}^2}{3600~\text{s}} = \frac{1}{3.6 \times 10^6}~\frac{\text{m}^3}{\text{s}}
\end{equation}

In municipal survey datasets where baseline peak runoff $Q_0(v)$ and baseline rainfall intensity $i_0(v)$ are cataloged, the contributing tributary area $A(v)$ can be analytically inverted via:
\begin{equation}
A(v) = \text{clip}\left(\frac{Q_0(v) \cdot 3.6 \times 10^6}{C(v) \cdot i_0(v)}, A_{\min}, A_{\max}\right)
\label{eqn:area_inversion}
\end{equation}
where $A_{\min} = 5{,}000~\text{m}^2$ and $A_{\max} = 50{,}000~\text{m}^2$ enforce physically plausible catchment boundaries for urban street segments, eliminating division-by-zero anomalies in survey records. For conduits lacking baseline runoff records, a standard urban tributary area of $20{,}000~\text{m}^2$ is adopted.

\section{Effective Capacity and Utilisation Formulation}
Every conduit $v$ is characterized by an unobstructed design conveyance capacity $Q_{\text{design}}(v)$ (in $\text{m}^3/\text{s}$), determined by cross-sectional geometry, bed slope, and roughness according to Manning's equation (Eqn.~\ref{eqn:manning_formula}). However, the accumulation of solid waste, silt, and urban sediment diminishes the effective cross-sectional flow area. Let $S(v) \in [0, 100]$ represent the percentage volumetric silt blockage. The effective available conveyance capacity $Q_{\text{avail}}(v)$ is formulated as:
\begin{equation}
Q_{\text{avail}}(v) = Q_{\text{design}}(v) \cdot \max\left(1.0 - \frac{S(v)}{100}, 0.05\right)
\label{eqn:avail_capacity}
\end{equation}
where a lower floor of $0.05$ (5\% residual capacity) prevents mathematical singularity when a channel is completely blocked.

The hydraulic capacity utilisation ratio $U(v, i)$ at rainfall intensity $i$ is defined as the ratio of incoming runoff to effective available conveyance capacity:
\begin{equation}
U(v, i) = \min\left(U_{\max}, \frac{Q(v, i)}{Q_{\text{avail}}(v)}\right)
\label{eqn:utilisation_def}
\end{equation}
where $U_{\max} = 3.0$ (300\% capacity utilisation) caps extreme ratios for stable numerical aggregation and visual map styling. 

Based on the calculated utilisation ratio, conduits are partitioned into three operational threat tiers:
\begin{equation}
T(v, i) = \begin{cases}
0~(\text{Low}), & \text{if } U(v, i) \le 0.80 \quad (\text{flow within safe design parameters}) \\
1~(\text{High}), & \text{if } 0.80 < U(v, i) \le 1.25 \quad (\text{near capacity; surcharge imminent}) \\
2~(\text{Critical}), & \text{if } U(v, i) > 1.25 \quad (\text{capacity breached; surface flooding active})
\end{cases}
\label{eqn:threat_tier}
\end{equation}

\section{Directed Acyclic Graph (DAG) Construction}
Urban stormwater flows strictly by gravity from higher topographic elevations toward lower discharge channels and water bodies. To model network interactions, CivicTwin constructs a Directed Acyclic Graph $G = (V, E)$, where vertices $V$ represent individual conduit segments ($|V| = 13{,}998$) and directed edges $E$ represent hydraulic connectivity ($|E| = 22{,}737$).

Constructing this topological graph from raw geospatial lines requires resolving junction connectivity:
\begin{enumerate}
    \item \textbf{Vertex Coordinate Extraction:} OpenStreetMap (OSM) drain ways meet at shared node junctions. Every polyline vertex coordinate $(lat, lon)$ is rounded to six decimal places ($\approx 0.11$~meter geographic precision) to serve as a topological junction key.
    \item \textbf{Elevation-Enforced Directionality:} Invert elevations $z(v)$ derived from the 30-meter SRTM DEM are assigned to all conduits. To guarantee the physical reality of downhill gravity flow and rigorously eliminate topological cycles, edges between conduits meeting at a shared junction are directed strictly from higher elevation to lower elevation:
    \begin{equation}
    (u, v) \in E \iff u \text{ and } v \text{ share a vertex, and } z(u) > z(v)
    \end{equation}
    Ties in elevation ($z(u) = z(v)$) are broken deterministically by internal segment index to ensure strict DAG properties.
    \item \textbf{Junction Edge Reduction:} If $k$ conduits meet at a shared junction point, instead of inserting all $\mathcal{O}(k^2)$ pairwise combinations, the conduits are sorted in descending elevation order:
    \begin{equation}
    u_1 \succ u_2 \succ \dots \succ u_k
    \end{equation}
    Directed edges are created between adjacent pairs $(u_j, u_{j+1})$ for $j = 1, \dots, k-1$. This chains the junction with exactly $k-1$ edges while preserving identical upstream and downstream reachability sets, minimizing graph edge density.
\end{enumerate}

\begin{algorithm}[h!]
\caption{Elevation-Enforced Drainage DAG Construction}
\label{alg:dag_construction}
\begin{algorithmic}[1]
\REQUIRE Set of conduits $V$ with polylines $P(v)$ and elevations $z(v)$
\ENSURE Directed Acyclic Graph $G = (V, E)$
\STATE Initialize directed graph $G \leftarrow (V, \emptyset)$
\STATE Initialize hash map $Junctions \leftarrow \emptyset$ \COMMENT{Map: vertex coordinate $\rightarrow$ list of conduit IDs}
\FORALL{$v \in V$}
    \FORALL{point $p \in P(v)$}
        \STATE $key \leftarrow (\text{round}(p.lat, 6), \text{round}(p.lon, 6))$
        \STATE Append $v$ to $Junctions[key]$
    \ENDFOR
\ENDFOR
\STATE Sort conduits $V$ by elevation descending: $z(u) > z(v)$, breaking ties by segment ID
\FORALL{$key, segs \in Junctions$}
    \STATE Remove duplicate segments: $segs \leftarrow \text{unique}(segs)$
    \IF{$|segs| \ge 2$}
        \STATE Sort $segs$ according to descending elevation rank
        \FOR{$j = 0$ \TO $|segs| - 2$}
            \STATE Add directed edge $(segs[j], segs[j+1])$ to $G$
        \ENDFOR
    \ENDIF
\ENDFOR
\STATE Assert that $G$ is a Directed Acyclic Graph (DAG)
\RETURN $G$
\end{algorithmic}
\end{algorithm}

\section{Bitset-Accelerated Ancestor Reachability}
A fundamental insight of CivicTwin is that clearing a desilted conduit does not merely benefit that individual channel; it restores hydraulic outflow for every upstream channel that discharges through it. We define the \textbf{Cascade Saved} metric, $C_{\text{saved}}(v)$, as the total number of distinct upstream ancestor conduits whose runoff passes through conduit $v$:
\begin{equation}
C_{\text{saved}}(v) = |\text{anc}(v)| = |\{u \in V \mid u \rightsquigarrow v \text{ in } G\}|
\label{eqn:cascade_saved}
\end{equation}

Computing transitive closure on a 14,000-node graph via recursive traversal is computationally expensive for interactive web applications. CivicTwin achieves instantaneous computation by executing topological sort reachability over bitsets. 

Let $\tau = (\tau_1, \tau_2, \dots, \tau_n)$ denote the topological sort of $G$, such that for every directed edge $(u, v) \in E$, $u$ appears before $v$ in $\tau$. Each node $v$ maintains a bitmask vector $M(v) \in \{0, 1\}^n$, where the $k$-th bit is 1 if node $k$ is an ancestor of $v$, and 0 otherwise. Processing nodes in topological order allows the ancestor mask of $v$ to be computed via bitwise logical OR ($\lor$) over its immediate predecessors:
\begin{equation}
M(v) = \bigvee_{p \in \text{pred}(v)} \left(M(p) \lor (1 \ll p)\right)
\label{eqn:bitset_or}
\end{equation}

The cardinality $|\text{anc}(v)|$ is subsequently obtained via hardware-accelerated population count (\texttt{popcount} / \texttt{bit\_count}), which counts the set bits in constant time. This reduces the asymptotic complexity to $\mathcal{O}(|V| \cdot \lceil |V| / 64 \rceil)$, executing across 13,998 nodes in under 20 milliseconds on standard CPU hardware.

\section{Systemic Network Risk Exposure Metric}
When a downstream conduit experiences hydraulic capacity exhaustion ($U > 1.0$), water cannot discharge freely. Backwater surcharge propagates upstream, submerging the outfall of connected conduits and raising their effective water levels. 

To quantify this phenomenon across the municipality, CivicTwin defines the \textbf{Systemic Network Risk Exposure}, $E_{\text{net}}(i)$, as the sum over all conduits of the worst capacity utilisation encountered at or downstream of that conduit:
\begin{equation}
E_{\text{net}}(i) = \sum_{v \in V} \max_{u \in \text{desc}(v) \cup \{v\}} \left(U(u, i) \times 100\right)
\label{eqn:network_risk_exposure}
\end{equation}
where $\text{desc}(v)$ denotes the set of all downstream descendants of $v$. 

Evaluating Eqn.~\ref{eqn:network_risk_exposure} is performed efficiently using dynamic programming in reverse topological order (from drainage outfalls up toward headwaters):
\begin{equation}
D(v) = \max\left(U(v, i) \times 100, \max_{s \in \text{succ}(v)} D(s)\right)
\label{eqn:downstream_dp}
\end{equation}
where $D(v)$ represents the worst downstream utilisation bottleneck affecting conduit $v$.

Under this formulation, clearing a downstream bottleneck simultaneously reduces $D(v)$ for that conduit and for every one of its upstream ancestors. Consequently, desilting a major downstream trunk conduit yields an exposure reduction proportional to its cascade reachability $C_{\text{saved}}$, mathematically explaining why cascade-aware prioritization substantially outperforms local severity-only desilting.
